\htmlhr
\chapterAndLabel{Delegation Checker}{delegation-checker}

The Delegation Checker verifies that a class that is written as a
\emph{delegator} really does delegate.  A delegator is a class that
implements each of its methods by calling the same method on some other
object, called the \emph{delegate}.  Delegation is an alternative to
inheritance:  rather than reusing a superclass's implementation, the
delegator forwards each call to an object that it wraps.

Here is an example of a delegator:

\begin{Verbatim}
class MyEnumeration<T> implements Enumeration<T> {
  @Delegate private Enumeration<T> e;

  public boolean hasMoreElements() {
    return e.hasMoreElements();
  }

  public T nextElement() {
    return e.nextElement();
  }
}
\end{Verbatim}

Delegation is worth verifying because a delegator inherits more than an
implementation:  it also inherits the delegate's specification.  For
example, suppose that \<Enumeration.hasMoreElements()> has a conditional
postcondition.  Then \<MyEnumeration.hasMoreElements()> satisfies that
postcondition only because its body is exactly a call to
\<e.hasMoreElements()>.  A typo, an extra side effect, or a forgotten
override silently invalidates the reasoning that another type-checker
performs.  The Delegation Checker detects such mistakes.

The Delegation Checker is not a type system:  it defines no type
qualifiers and no subtyping relationship.  It is nonetheless a checker, and
it is run in the same way as any other checker.


\sectionAndLabel{Running the Delegation Checker}{delegation-checker-running}

An example invocation is

\begin{Verbatim}
javac -processor org.checkerframework.common.delegation.DelegationChecker MyFile.java
\end{Verbatim}

The Delegation Checker's purpose is to justify reasoning that
\emph{other} checkers perform, so it is most useful when run together with
the type-checkers whose specifications the delegator depends on:

\begin{Verbatim}
javac -processor MapKeyChecker,DelegationChecker MyFile.java
\end{Verbatim}


\sectionAndLabel{Annotations}{delegation-checker-annotations}

The Delegation Checker uses the following declaration annotations:

\begin{description}
\item[\refqualclass{common/delegation/qual}{Delegate}]
  is written on a field.  It indicates that the field is the delegate of
  the class that declares it:  the enclosing class's methods are
  implemented by calling the same method on the field.  A class may have at
  most one \<@Delegate> field.
\item[\refqualclass{common/delegation/qual}{DelegatorMustOverride}]
  is written on a method \<m> of a class \<C>.  It indicates that any
  delegator that is a subtype of \<C> must override \<m>.  Write it on a
  method whose specification (such as a precondition, a postcondition, or a
  conditional postcondition) would not be satisfied by an inherited
  implementation.
\end{description}


\sectionAndLabel{What the Delegation Checker guarantees}{delegation-checker-guarantees}

Within a class that declares a field annotated \<@Delegate>, the Delegation
Checker enforces the following rules.

\begin{itemize}
\item
  The class declares at most one field annotated \<@Delegate>.  Otherwise,
  the checker issues a \<multiple.delegate.annotations> error.
\item
  The body of every method that overrides (or implements) a supertype
  method consists of a single statement, which is
  \begin{itemize}
  \item
    a call to the same method on the delegate field, passing the
    overriding method's formal parameters, in order, as the arguments; or
  \item
    \<throw new UnsupportedOperationException(...)>.  A delegator may
    refuse to support an operation rather than delegating it.
  \end{itemize}
  Otherwise, the checker issues an \<invalid.delegate> warning.
\item
  The class overrides every method of a supertype that is annotated
  \<@DelegatorMustOverride>.  Otherwise, the checker issues a
  \<delegate.override> warning that lists the methods that were not
  overridden.
\end{itemize}

The Delegation Checker does not require a delegator to override
\emph{every} method of its supertype.  Requiring that would be onerous, and
it is unnecessary for a method whose inherited implementation is already
correct for the delegator.  Use \<@DelegatorMustOverride> to mark the
methods for which an inherited implementation is not acceptable.

A class with no field annotated \<@Delegate> is not a delegator, so the
Delegation Checker does not check it.


\sectionAndLabel{Example}{delegation-checker-example}

\begin{Verbatim}
import java.util.IdentityHashMap;
import java.util.Map;
import org.checkerframework.common.delegation.qual.Delegate;

public class MyMap<K, V> extends IdentityHashMap<K, V> {

  @Delegate private final IdentityHashMap<K, V> map;

  MyMap(IdentityHashMap<K, V> map) {
    this.map = map;
  }

  @Override
  public int size() {
    return map.size();     // OK: delegates
  }

  @Override
  public V get(Object key) {
    return map.get(key);   // OK: delegates, passing the formal parameter
  }

  @Override
  public boolean isEmpty() {
    // Warning: the body is not a call to map.isEmpty()
    return map.size() == 0;
  }

  @Override
  public void putAll(Map<? extends K, ? extends V> m) {
    // OK: a delegator may refuse to support an operation
    throw new UnsupportedOperationException();
  }
}
\end{Verbatim}


% LocalWords:  delegator delegators MyEnumeration hasMoreElements nextElement
% LocalWords:  DelegatorMustOverride IdentityHashMap MyMap
% LocalWords:  UnsupportedOperationException putAll MapKeyChecker
% LocalWords:  DelegationChecker
